\section{Impact Modeling}

\subsection{Model Overview}

We modeled wildlife population dynamics under three major external pressures: wildfire, poaching, and tourism/urbanization. These impacts were chosen because they represent distinct types of ecological stress. Wildfire is a natural environmental disturbance that intensifies over time. Poaching represents a human-driven process motivated by economic incentives. Tourism/urbanization introduces continuous, infrastructure-dependent habitat disruption.

Our model captures how each of these processes interacts with wildlife populations over time through two mechanisms: direct mortality (death caused by the disturbance) and habitat degradation (reduction in ecosystem carrying capacity via resource destruction). By combining ecological processes (e.g.\ fire spread, habitat recovery) with human components (e.g.\ poacher decision-making, tourism pressure), our model serves as a dynamic representation of wildlife systems in Etosha National Park under various impacts.

\subsection{List of Variables}

\begin{table}[H]
\centering
\[
\begin{array}{lll}
\hline
\textbf{Symbol} & \textbf{Description} & \textbf{Units} \\
\hline
t & \text{Time} & \text{days} \\
s & \text{Species index} & - \\
N_s(t) & \text{Population of species } s & \text{individuals} \\
R_m(t),\,R_m^0 & \text{Amount of resource } m \text{ (current, baseline)} & \text{resource units} \\
r_{m,s} & \text{Per-capita requirement of resource } m \text{ for species } s & \text{resource units/individual} \\
K_s(t) & \text{Carrying capacity of species } s & \text{individuals} \\
r_h^{(m)} & \text{Natural recovery rate of resource } m & \text{day}^{-1} \\
\mathcal{M}_s(t) & \text{Total threat-induced mortality rate for species } s & \text{individuals/day} \\
F(t) & \text{Wildfire intensity} & - \\
P(t) & \text{Poaching effort} & - \\
T(t) & \text{Tourism disturbance level} & - \\
u(t) & \text{Fire suppression effort (from optimization)} & - \\
\hline
\end{array}
\]
\caption{General model variables}
\label{tab:variables}
\end{table}

\subsection{Base Population Dynamics and Carrying Capacity}

\subsubsection{Carrying Capacity}

All species share a common resource pool. Let $R_m(t)$ denote the available amount of resource $m$ (e.g.\ water, forage, space) at time $t$, with baseline level $R_m^0$. Each species $s$ requires $r_{m,s}$ units of resource $m$ per individual. The carrying capacity of species $s$ is then determined by Liebig's law of the minimum:
\[
K_s(t) = \min_m \frac{R_m(t)}{r_{m,s}}.
\]
This ensures that the most limiting resource constrains the population. Resources evolve over time according to the combined effects of natural recovery, wildfire damage, and tourism degradation (see Sections~\ref{sec:wildfire} and~\ref{sec:tourism}):
\[
\frac{dR_m}{dt} = r_h^{(m)}\bigl(R_m^0 - R_m\bigr) - \lambda_h^{(m)} F(t)\,R_m - \lambda_{\text{tour}}^{(m)} T(t)\,R_m.
\]
The first term drives recovery toward baseline, while the latter two terms capture resource destruction from fire and tourism respectively. Recovery rates are modulated by rainfall seasonality: water recovery scales directly with a normalized rainfall index, forage recovery lags rainfall by $\approx 30$~days, and space recovery is rainfall-independent. A dry-season floor of $0.3$ reflects Etosha's permanent groundwater-fed waterholes \cite{craig2015etosha}.

\begin{center}
\begin{tabular}{l c c c l}
\hline
\textbf{Resource} & $R_m^0$ & $r_h^{(m)}$ & \textbf{Binding species} & \textbf{Source / Rationale} \\
\hline
Water & 1000 & 0.04\,day$^{-1}$ & elephant ($r = 0.15$) & permanent springs recover fast \\
Forage & 1000 & 0.05\,day$^{-1}$ & elephant ($r = 0.12$) & grassland recovers in wet season \\
Space & 1000 & 0.01\,day$^{-1}$ & lion ($r = 0.12$) & habitat structure slow to recover \\
\hline
\end{tabular}
\end{center}

Initial populations: $N_{\text{zebra}}(0) = 13{,}000$, $N_{\text{lion}}(0) = 500$, $N_{\text{elephant}}(0) = 2{,}500$, $N_{\text{rhino}}(0) = 70$ \cite{craig2015etosha}.

\subsubsection{Predator--Prey Dynamics (Zebra and Lion)}

Zebras and lions are coupled through a Lotka--Volterra predator--prey system. Writing $N_z$ and $N_l$ for zebra and lion populations:
\begin{align}
\frac{dN_z}{dt} &= \lambda_1 N_z\!\left(1 - \frac{N_z}{K_z(t)}\right) - \gamma\, N_z\, N_l - \mathcal{M}_z(t), \label{eq:zebra} \\[6pt]
\frac{dN_l}{dt} &= \lambda_2\, N_z\, N_l - \delta\, N_l - \zeta\, N_l^2 - \mathcal{M}_l(t), \label{eq:lion}
\end{align}
where $\mathcal{M}_s(t)$ denotes the total threat-induced mortality rate for species $s$ (defined below). Zebra growth is logistic with a time-varying carrying capacity and is reduced by predation. Lion growth depends on prey availability and is checked by natural death $\delta$ and intraspecific competition $\zeta$.

\begin{center}
\begin{tabular}{l l c l}
\hline
\textbf{Parameter} & \textbf{Definition} & \textbf{Value} & \textbf{Source} \\
\hline
$\lambda_1$ & zebra intrinsic growth rate & $0.0015\;\text{day}^{-1}$ & $\approx 55\%$/yr \cite{gaillard2000dynamics} \\
$\gamma$ & predation rate & $1 \times 10^{-6}$ & $\approx 6.5$ zebra/day at equilibrium \cite{hayward2005prey} \\
$\lambda_2$ & lion conversion efficiency & 0.14 & fitted to observed lion growth \cite{stander1991demography} \\
$\delta$ & lion natural death rate & $3 \times 10^{-4}\;\text{day}^{-1}$ & $\approx 9$\,yr lifespan \cite{barthold2016lion} \\
$\zeta$ & lion intraspecific competition & $2 \times 10^{-6}$ & density-dependent regulation \cite{kissui2004lions} \\
$r_{\text{elephant}}$ & elephant intrinsic growth & $3 \times 10^{-4}\;\text{day}^{-1}$ & $\approx 4\%$/yr \cite{calef1988elephant} \\
$r_{\text{rhino}}$ & black rhino intrinsic growth & $4 \times 10^{-4}\;\text{day}^{-1}$ & $\approx 5\%$/yr \cite{emslie1999rhino} \\
\hline
\end{tabular}
\end{center}

\subsubsection{Logistic Growth (Other Species)}

All other species (rhinos, elephants, leopards, cheetahs, birds, plants) follow logistic growth:
\[
\frac{dN_s}{dt} = r_s\, N_s\!\left(1 - \frac{N_s}{K_s(t)}\right) - \mathcal{M}_s(t).
\]

\subsubsection{Threat Mortality}

The total threat-induced mortality for each species is the sum of wildfire, poaching, and tourism contributions:
\[
\mathcal{M}_s(t) = \underbrace{\mu_s(t)\,F(t)\,N_s}_{\text{wildfire}} + \underbrace{\alpha_s(t)\,P(t)\,N_s}_{\text{poaching}} + \underbrace{\beta_s\,T(t)\,N_s}_{\text{tourism}}.
\]
Each component is detailed in the sections that follow.

\subsection{Wildfire}
\label{sec:wildfire}

Wildfire is modeled as an ecological disturbance driven by fuel availability, dryness, and human suppression. Fire causes direct mortality and degrades resources, which in turn reduces carrying capacity. Resource recovery after fire is captured by the $r_h^{(m)}(R_m^0 - R_m)$ term in the resource ODE.

\subsubsection{Wildfire Model Parameters}

\begin{center}
\begin{tabular}{l l c l}
\hline
\textbf{Parameter} & \textbf{Definition} & \textbf{Value} & \textbf{Source / Rationale} \\
\hline
$\gamma_0$ & fire spread rate & 0.6\,day$^{-1}$ & moderate savanna spread \\
$d_0,\,d_1$ & dryness baseline, amplitude & 0.2, 0.8 & strong Etosha seasonal variation \\
$t_{\text{dry}}$ & peak dry timing & 180\,days & mid-year dry season peak \\
$c_0,\,c_1$ & passive, active suppression & 0.1, 0.5 & active control halves spread \\
$\mu_{\text{zebra}}^{\text{base}}$ & zebra fire mortality & $1.5 \times 10^{-4}$ & high mobility reduces risk \\
$\mu_{\text{lion}}^{\text{base}}$ & lion fire mortality & $5 \times 10^{-4}$ & moderate agility \\
$\mu_{\text{elephant}}^{\text{base}}$ & elephant fire mortality & $2 \times 10^{-4}$ & large body, some mobility \\
$\mu_{\text{rhino}}^{\text{base}}$ & black rhino fire mortality & $3 \times 10^{-4}$ & low mobility, small range \cite{emslie1999rhino} \\
$\lambda_h^{(\text{water})}$ & fire damage to water & 0.01 & fire doesn't drain waterholes \\
$\lambda_h^{(\text{forage})}$ & fire damage to forage & 0.02 & grassland burns but recovers \\
$\lambda_h^{(\text{space})}$ & fire damage to space & 0.01 & habitat structure less affected \\
\hline
\end{tabular}
\end{center}

\subsubsection{Wildfire Dynamics}

Wildfire intensity $F(t) \in [0,1]$ evolves as:
\[
\frac{dF}{dt} = \gamma_f\bigl(d(t), F(t)\bigr)\,F(t) - c(t)\,F(t),
\]
where the spread rate is
\[
\gamma_f(d,F) = \gamma_0\,d(t)\,(1-F).
\]
This captures three ecological facts: fire spreads faster in dry conditions, growth slows as fuel is consumed (the $1-F$ saturation), and spread is proportional to current fire extent.

Dryness follows a seasonal cycle:
\[
d(t) = d_0 + d_1\max\!\left(0,\,\sin\!\left(\frac{2\pi(t - t_{\text{dry}})}{365}\right)\right),
\]
where the $\max$ ensures only the dry phase contributes to fire risk. Fire suppression is
\[
c(t) = c_0 + c_1\,u(t),
\]
where $u(t)$ represents human intervention such as firefighting or controlled burns.

\subsubsection{Wildfire Mortality}

Direct fire mortality for species $s$ contributes to $\mathcal{M}_s$ as:
\[
\mu_s(t)\,F(t)\,N_s(t),
\]
where the mortality coefficient varies seasonally:
\[
\mu_s(t) = \mu_s^{\text{base}}\left[1 + 0.3\max\!\left(0,\,\sin\!\left(\frac{2\pi(t - t_{\text{dry}})}{365}\right)\right)\right].
\]
Fires during peak dry periods are more severe, producing higher mortality rates.

\subsubsection{Resource Damage}

Wildfire degrades resources through the $\lambda_h^{(m)} F(t)\,R_m$ term in the resource ODE (Section~3.3.1). The parameter $\lambda_h^{(m)}$ is resource-specific: grassland resources recover quickly ($r_h^{\text{grass}} = 0.03$) while forest resources regenerate slowly ($r_h^{\text{forest}} = 0.01$). As resources decline, $K_s(t)$ falls, naturally suppressing population growth.

\subsection{Poaching}

Poaching is modeled as an adaptive human activity driven by economic incentives and constrained by enforcement. We treat poachers as rational agents who adjust effort based on expected profit.

\subsubsection{Species-Specific Poaching Parameters}

\begin{center}
\begin{tabular}{lccl}
\hline
\textbf{Species} & $\alpha_s^{\text{base}}$ & $w_s$ & \textbf{Source / Rationale} \\
\hline
Black rhino & $3 \times 10^{-4}$ & 1.5 & highest horn value; critically endangered \cite{emslie1999rhino} \\
Elephant & $3 \times 10^{-4}$ & 1.0 & strong ivory demand \cite{wittemyer2014ivory} \\
Lion & $1 \times 10^{-3}$ & 0.7 & trophy hunting demand \cite{trinkel2013lions} \\
Plains zebra & $1 \times 10^{-4}$ & 0.08 & minimal targeting \cite{gaillard2000dynamics} \\
\hline
\end{tabular}
\end{center}

Additional poaching dynamics parameters: effort adjustment rate $\gamma_p = 0.001$\,day$^{-1}$, revenue coefficient $r_p = 5 \times 10^{-5}$, opportunity cost $\Omega = 0.08$, risk sensitivity $\sigma_p = 0.5$, baseline enforcement $\Lambda_0 = 0.1$, and tourism enforcement boost $\lambda = 0.3$.

\subsubsection{Poaching Dynamics}

Poaching effort $P(t)$ evolves as:
\[
\frac{dP}{dt} = \gamma_p\left[r_p\,\mathcal{N}(t) - \bigl(\Omega + \sigma_p(\Lambda_0 + \lambda\, T(t))\bigr)\right]P(t),
\]
where $\mathcal{N}(t) = \sum_s w_s\, N_s(t)$ is the value-weighted prey pool. Poaching effort grows when expected revenue exceeds the sum of opportunity cost $\Omega$ and risk-adjusted enforcement $\sigma_p(\Lambda_0 + \lambda\, T(t))$. The tourism term means that increased human presence raises enforcement and discourages poaching.

\subsubsection{Poaching Mortality}

The poaching contribution to $\mathcal{M}_s$ is:
\[
\alpha_s(t)\,P(t)\,N_s(t),
\]
with seasonal variation:
\[
\alpha_s(t) = \alpha_s^{\text{base}}\left[1 - 0.3\max\!\left(0,\,\sin\!\left(\frac{2\pi(t - t_{\text{wet}})}{365}\right)\right)\right].
\]
During wetter periods, animals are more dispersed and harder to track, reducing poaching efficiency.

\subsection{Tourism and Urbanization}
\label{sec:tourism}

Tourism and urbanization introduce persistent human disturbance through repeated human presence and infrastructure development. Unlike wildfire and poaching, tourism pressure is continuous rather than episodic.

\subsubsection{Tourism and Urbanization Parameters}

\begin{center}
\begin{tabular}{l l c l}
\hline
\textbf{Parameter} & \textbf{Definition} & \textbf{Value} & \textbf{Source / Rationale} \\
\hline
$w_v,\,w_w,\,w_p$ & disturbance weights & 0.5, 0.3, 0.2 & visitors dominate disturbance \\
$V_{\max}$ & max daily visitors & 1000 & sustainable park capacity \\
$W_{\max}$ & max waterhole crowding & 500 & limited water availability \\
$P_{\text{infra}}$ & permanent infrastructure & 0.15 & normalized lodge/road footprint \\
$\beta_{\text{zebra}}$ & zebra tourism mortality & $1 \times 10^{-4}$ & tolerant of human presence \cite{gaillard2000dynamics} \\
$\beta_{\text{lion}}$ & lion tourism mortality & $2 \times 10^{-4}$ & avoids but tolerates humans \cite{trinkel2013lions} \\
$\beta_{\text{elephant}}$ & elephant tourism mortality & $2 \times 10^{-4}$ & large range increases contact \cite{calef1988elephant} \\
$\beta_{\text{rhino}}$ & black rhino tourism mortality & $3 \times 10^{-4}$ & highly sensitive to disturbance \cite{emslie1999rhino} \\
$\lambda_{\text{tour}}^{(\text{water})}$ & tourism water damage & 0.005 & infrastructure affects water \\
$\lambda_{\text{tour}}^{(\text{forage})}$ & tourism forage damage & 0.002 & trampling and clearing \\
$\lambda_{\text{tour}}^{(\text{space})}$ & tourism space damage & 0.008 & roads and lodges reduce habitat \\
\hline
\end{tabular}
\end{center}

\subsubsection{Tourism Disturbance}

The tourism disturbance index is:
\[
T(t) = w_v\frac{V(t)}{V_{\max}} + w_w\frac{W(t)}{W_{\max}} + w_p\, P_{\text{infra}}.
\]
This aggregates visitor volume, water usage, and infrastructure development into a normalized disturbance level.

\subsubsection{Tourism Mortality}

Tourism contributes to $\mathcal{M}_s$ as:
\[
\beta_s\,T(t)\,N_s(t).
\]
The coefficient $\beta_s$ captures species-specific sensitivity to human disturbance; species that are more easily stressed by human presence experience higher effective mortality.

\subsubsection{Resource Damage}

Tourism degrades resources through the $\lambda_{\text{tour}}^{(m)} T(t)\,R_m$ term in the resource ODE (Section~3.3.1). Infrastructure development causes persistent damage (reflected in the constant $P_{\text{infra}}$ component of $T(t)$), while visitor-driven disturbance fluctuates. As with wildfire, resource loss reduces $K_s(t)$ and naturally suppresses population growth.

\subsection{Seasonal Application of Model over Time}

To understand the long-term sustainability of our model, we hope to characterize the oscillating temporal shifts in resource availability, how those changes affect migration patterns, and to finally understand how our model might respond to new predicted risk landscapes and animal distributions. 

Biannual oscillations in Etosha weather trends are shaped by a "dry" and "wet" season, named for the parallel patterns seen in rainfall metrics. Currently, our model operates under data taken from a dry winter, where animals and risks are concentrated around water holes, a consequence of water scarcity turning the salt pan and its surroundings on the outskirts of the park into an arid desert. Moreover, we note that space and vegetation, two other limiting resources, often only become limiting as a result of altered water availability. These concentrated distributions expand outward over time, as wildlife migrates in search of less competition across a lush bushland laden with freshly-filled water holes. However, these changes aren't permanent, as the same tropical atmospheric circulation that created the abundance saps moisture from the region, once again settling into a dry season. 

Thus, as the seasonal changes are primarily a function of atmospheric change, with an emphasis on rainfall or humidity metrics, we use the RAINFALL INDEX to fit a function and model the temporal expansion of water holes over time. 

INSERT MULTI-PANEL FIGURE

In order for rangers to better surveil wildlife through the thick foliage during the peaks of wet season, we  provide a model that predicts the spatial migration of wildlife as a function of increased rainfall. To do so, we would apply the periodic function to a spatial model of wildlife distribution for each species, as taken from tracking data and recorded sightings around the central water-holes, and couple the resulting expanded range with a maximum entropy function to map populations to grids. For more fine-tuned, regional analysis, we could create a new displacement function that is modeled after spatial diffusion mechanics wherein populations slowly move towards surrounding cells with a high carrying capacity and low populations in effort to meet a greater global attractor. This could be done by conditionally adding a term to the target cells and subtracting the same term from the original cell. This term can be modeled by the simple partial differential equation (PDE): $$\frac{\delta u_s}{\delta t} = D_s\nabla^2u_s$$ where $u_s$ is population density of a given species, $\nabla^2$ is the Laplace operator, and $D_s$ is the diffusion coefficient, which may be derived from the above rainfall indices. 



Interestingly, we notice that each of our threats should be neutralized as a consequence of the combination of landscape changes and animal density. As habitat fragmentation naturally decreases when resources are once again scattered throughout the region, the threats of wildfires and tourism will now naturally not generate as many deaths, as animal density is reduced over time. Additionally, the humid environment and lush greenery during the peaks of wet season ensures that wildfires have essentially no chance of spreading, or even sparking. Finally, while poaching may seem more difficult to trap, prevent, and surveil as a result of the expansive distribution of animals, we note that poachers themselves struggle to reliably hunt and capture elephants and black rhinos during wet season among the intense foliage for many of the same reasons. 

\textbf{EITHER WE HYPOTHESIZE HERE or we conclude} that our model of surveillance and ranger distribution can be kept the same to even being relaxed during this time period, as dry season is clearly the worst-case scenario. 

\url{https://www.safaribookings.com/etosha/climate}
